\documentclass[12pt,a4paper]{article}

\usepackage[UTF8,fontset=fandol]{ctex}
\usepackage{amsmath,amssymb,amsthm,mathtools}
\usepackage{geometry}
\geometry{margin=2.5cm}
\usepackage{hyperref}
\hypersetup{colorlinks=true,linkcolor=blue!70!black,urlcolor=blue!60!black}
\usepackage{enumitem}
\setlist{nosep,leftmargin=*}
\usepackage{tikz}
\usetikzlibrary{arrows.meta,positioning}

\theoremstyle{definition}
\newtheorem{definition}{定义}[section]
\newtheorem{example}{例}[section]
\newtheorem{exercise}{练习题}

\theoremstyle{plain}
\newtheorem{problem}{问题}[section]
\newtheorem{proposition}[problem]{命题}

\theoremstyle{remark}
\newtheorem*{remark}{注记}
\newtheorem*{insight}{关键 Insight}
\newtheorem*{pitfall}{典型错因}

\title{立方体上屏：一条解析几何管线}
\author{方块黎明数学讲义}
\date{2026-11-16}

\begin{document}
\maketitle
\tableofcontents
\bigskip

\noindent
这份讲义只解释当前程序里那一个立方体。约定一次说清，后面所有推导都沿用它。

\begin{itemize}
\item 右手系。$+X$ 向右，$+Y$ 向上。相机放在 $+Z$ 一侧看向原点，因此视线沿着 $-Z$，相机前方的点满足 $z<0$。
\item 点是列向量。若干变换连乘时，靠近点的矩阵先发生：$p_{\mathrm{clip}}=PVM\,p$。
\item 立方体中心在原点，边长为 $1$，八个角的坐标分量都是 $\pm 1/2$。
\item 世界由种子 $20261116$ 生成。$i,k\in[-10,10]$ 的每一列有 $4$ 块实心方块，顶面高度由平滑后的散列决定。人从方块 $(2,3)$ 的上空出生，视线对准原点。偏航、俯仰、行走、重力和碰撞仍按前面的约定。竖直视野 $50^\circ$，近平面距离 $n=0.1$，远平面距离 $f=100$。
\end{itemize}

\section{动机}

画面上的立方体要先回答四个几何问题：立方体由哪些点构成；这些点怎样转起来；眼睛在哪里、朝哪里；近处的同样长度为什么看起来更长。前三个是空间里的刚体与坐标系，第四个是透视。

透视这条要求很老。Brunelleschi 约在 1415 年用灭点做了透视演示，Alberti 在 1435 年的 \emph{De pictura} 里把看见的范围写成一只视觉金字塔：眼睛是顶点，画面是截这条金字塔的一块平面。近处一段和远处一段在画面上切出的长度不同。这就是近大远小的几何来源，当时还没有矩阵。

矩阵要等齐次坐标。M\"obius 在 1827 年的 \emph{Der barycentrische Calcul} 里用重心坐标处理射影变换。平移在 $\mathbb{R}^3$ 里不是线性映射，加进第四个坐标之后，平移、旋转、透视可以连成一串 $4\times 4$ 矩阵。计算机图形学把这条链用成标准管线，早期的三维透视机器图见于 Roberts 1963 年的博士论文 \emph{Machine Perception of Three-Dimensional Solids}。OpenGL 固定的是这条链的一个具体约定：相机看向 $-Z$，裁剪之后的可见范围是 $[-1,1]^3$。

当前程序把这条链停在顶点的屏幕坐标和深度次序。三角形内部填什么颜色，由显卡按顶点颜色插值完成。那一步仍是几何，但这份讲义的推导停在顶点。

\section{形式定义}

\begin{definition}[点、向量、齐次坐标]
空间中的点写成 $p=(x,y,z)$。向量也写成三个数，但平移只作用于点：点加向量仍是点，向量加向量仍是向量。齐次坐标把点写成
\[
\tilde p=\begin{pmatrix}x\\y\\z\\1\end{pmatrix},
\]
把向量写成第四个分量为 $0$ 的列。两个齐次列 $(x,y,z,w)$ 与 $(\lambda x,\lambda y,\lambda z,\lambda w)$ 在 $w\neq 0$、$\lambda\neq 0$ 时表示同一个点，普通坐标是 $(x/w,y/w,z/w)$。
\end{definition}

\begin{definition}[模型、观察、投影]
设立方体局部坐标中的点为 $p_{\mathrm{local}}$。三张 $4\times 4$ 矩阵依次为
\begin{align*}
M&\colon \text{局部坐标}\to\text{世界坐标},\\
V&\colon \text{世界坐标}\to\text{相机坐标},\\
P&\colon \text{相机坐标}\to\text{裁剪坐标}.
\end{align*}
裁剪坐标 $p_{\mathrm{clip}}=(x_c,y_c,z_c,w_c)$ 再做一次除法
\[
(x_n,y_n,z_n)=\Bigl(\frac{x_c}{w_c},\frac{y_c}{w_c},\frac{z_c}{w_c}\Bigr),
\]
得到标准化设备坐标。可见范围是三个分量都落在 $[-1,1]$，并且 $w_c>0$。
\end{definition}

\begin{definition}[视口]
帧缓冲宽 $W$ 像素、高 $H$ 像素，原点在左下。标准化设备坐标变到像素坐标为
\begin{equation}
x_{\mathrm{pix}}=\frac{x_n+1}{2}W,\qquad
y_{\mathrm{pix}}=\frac{y_n+1}{2}H.
\label{eq:viewport}
\end{equation}
\end{definition}

\begin{remark}
定义里真正产生近大远小的是除以 $w_c$ 那一步。$P$ 本身是线性的。把除法提前吸进 $P$，这张矩阵就不再是线性映射。
\end{remark}

\section{多视角表征}

\subsection{Mental Picture：脑海图像 / 草稿纸 / 符号层}

\textbf{脑海图像}：原点上立着一个边长为 $1$ 的玻璃立方体。眼睛在 $(1.5,1.2,2.5)$，视线钻进立方体中心。立方体绕竖直轴转时，八个角在空间里画圆，画面上的位置却不是匀速圆周，因为还有透视。

\textbf{草稿纸图像}：先画五格横排
\[
\text{局部}\xrightarrow{M}\text{世界}\xrightarrow{V}\text{相机}\xrightarrow{P}\text{裁剪}\xrightarrow{\div w}\text{屏幕}.
\]
再单独画 $xz$ 平面上的单位圆，标出 $+Z$ 经正角度转向 $+X$。最后画一张侧视的视锥：眼睛在原点，近平面在 $z=-n$，远平面在 $z=-f$，要保留的是两平面之间的那一段。

\textbf{符号层映射}：玻璃立方体的角是 $p_{\mathrm{local}}$。转动是 $M=R_y(\omega t)$。眼睛、目标、上方向三个箭头决定 $V$ 的三行。视锥的张角和远近决定 $P$。除法之后的 $(x_n,y_n)$ 才是草稿纸上那张画面里的位置，$z_n$ 是谁挡住谁。

\textbf{图像纠错}：常见的错图是把眼睛放在原点、让前方变成 $z>0$，同时仍使用下面这张最后一行是 $(0,0,-1,0)$ 的 $P$。那张 $P$ 默认前方是 $z<0$。两套图叠在一起时，$w_c$ 变成负数，立方体整颗落在可见范围外面。正确的替代图是：侧视图里，眼睛看向纸面的左方，近处那条竖线在 $z=-0.1$，立方体在更左的 $z\approx -3$ 附近。

\subsection{视角 A：几何}

立方体是八个点的凸包。一个面是其中四个共面的点。四个点不构成三角形，所以沿一条对角线切开，一个面变成两个三角形。三角形的朝向用叉积读出来：若 $A,B,C$ 是顶点顺序，则 $(B-A)\times(C-A)$ 指向「从这个方向看，三角形是逆时针」的那一侧。程序要求这个方向等于面的外法向，于是从立方体内部看每个面都是顺时针，背面剔除会丢掉它们。

\subsection{视角 B：线性代数}

旋转、换相机、把视锥压进立方体 $[-1,1]^3$，这三件事在齐次坐标里都是 $4\times 4$ 矩阵。连乘的顺序就是几何作用的反序书写：先旋转的 $M$ 写在最右边。透视和其他两步不同，它的线性部分故意把深度写进 $w$，真正的缩放留到除法。

\begin{example}[画面中心]
世界原点是立方体中心，也是相机的目标。它在相机坐标里位于视线轴上，所以 $x_c=y_c=0$。除法之后仍有 $x_n=y_n=0$。按 \eqref{eq:viewport}，在 $960\times 540$ 的帧缓冲上它落在像素 $(480,270)$，也就是画面中心。深度并不在中心：下面的数值例子给出 $z_n\approx 0.9385$，靠近远平面一侧的 $1$，但仍然可见。
\end{example}

\section{演化脉络}

\textbf{前置问题}：$3\times 3$ 矩阵能旋转、能缩放，不能把立方体从原点挪到别处，也不能把视锥变成屏幕上的矩形。平移会把原点送到非零向量，线性映射必须把原点送到原点。

\textbf{关键洞察}：M\"obius 的齐次坐标多出一个权重 $w$。点的权重取 $1$ 之后，平移变成第四列；透视则把权重改写成深度。图形管线把这两次「线性代数做不到的事」都推迟到第四个坐标上。

\textbf{后续影响}：同一串矩阵可以每帧更换。立方体转动只改 $M$，窗口拉宽只改 $P$ 里的宽高比，眼睛移动只改 $V$。显卡不需要知道这些矩阵的几何来历，它只做乘法和一次除法。

\begin{insight}
近大远小不在 $P$ 的乘法里完成。$P$ 只把深度抄进 $w_c$。所有点仍然被线性映射。近处显得大，是因为随后用一个更小的 $w_c$ 去除。透视的非线性全部集中在这一次除法。
\end{insight}

\section{卡点分析}

\textbf{老板原话}：``是不是本身是一个三维的数学解析几何问题啊''

\textbf{卡点位置}：已经看出立方体、旋转和相机是解析几何，但还没有把「齐次线性映射」和「除以 $w$」分成前后两步。八个角与 $36$ 个顶点的差别也还叠在同一句话里。

\textbf{卡点类型}：多视角切换。几何对象、矩阵、屏幕像素是三层，少一层都会把后面的公式代错地方。

\textbf{突破方式}：草稿纸上先只画五个坐标空间的箭头。每一张矩阵只回答一个箭头。除法单独成一步，不并进 $P$ 的口头描述里。顶点个数放到叉积和外法向之后再数。

\section{问题与推导}

\subsection{八个角为什么写成 36 个顶点}

\begin{problem}
立方体有 $8$ 个角。程序里的顶点数组有 $36$ 个顶点。这两个数各自在数什么？
\end{problem}

几何上，角是位置。一个角同时属于三个面。这一章每个面有自己的颜色，颜色是顶点属性，不是位置的属性。同一个位置若要带三种颜色，就要写成三条顶点记录。

一个面先有 $4$ 个角。沿对角线切成 $2$ 个三角形后，若不另做索引，每个三角形把 $3$ 个顶点完整写进数组。一个面占
\[
2\times 3=6
\]
条记录。六个面占
\[
6\times 6=36
\]
条记录。朝 $+Z$ 的面四个角是
\begin{align*}
A&=(-1/2,-1/2,1/2),&
B&=(1/2,-1/2,1/2),\\
C&=(1/2,1/2,1/2),&
D&=(-1/2,1/2,1/2).
\end{align*}
两条三角形是 $A\to B\to C$ 与 $A\to C\to D$。$A$ 和 $C$ 各出现两次。

检验外法向。对第一条，
\[
B-A=(1,0,0),\qquad C-A=(1,1,0),
\]
\[
(B-A)\times(C-A)
=\begin{vmatrix}
\mathbf i&\mathbf j&\mathbf k\\
1&0&0\\
1&1&0
\end{vmatrix}
=(0,0,1).
\]
这正是 $+Z$ 面的外法向。另外五个面用同样的算法排顺序，叉积分别等于 $(\pm 1,0,0)$、$(0,\pm 1,0)$、$(0,0,\pm 1)$ 中向外的那一个。

\subsection{平移为什么要第四个坐标}

\begin{problem}
证明：平移 $T(p)=p+t$（$t\neq 0$）不是 $\mathbb{R}^3$ 上的线性映射，但在齐次坐标下由一张 $4\times 4$ 矩阵表示。
\end{problem}

线性映射必须固定原点。$T(0)=t\neq 0$，所以 $T$ 不是线性的。也可以看加法：
\[
T(p+q)=p+q+t,\qquad T(p)+T(q)=p+q+2t.
\]
二者不相等。

把点写成 $(x,y,z,1)$。定义
\[
\tilde T=
\begin{pmatrix}
1&0&0&t_x\\
0&1&0&t_y\\
0&0&1&t_z\\
0&0&0&1
\end{pmatrix}.
\]
则
\[
\tilde T\begin{pmatrix}x\\y\\z\\1\end{pmatrix}
=\begin{pmatrix}x+t_x\\ y+t_y\\ z+t_z\\ 1\end{pmatrix}.
\]
第四个分量仍是 $1$，前三个分量正是 $T(p)$。$\tilde T$ 对四个坐标是线性的。向量的第四个分量是 $0$，乘出来平移列不起作用，所以向量不被平移。这就是点和向量必须分开写的原因。

这一章的立方体还停在原点，模型矩阵里的平移是单位矩阵。齐次坐标仍然必要，因为下一步的相机平移和透视都要用第四列或第四行。

\subsection{绕竖直轴的旋转}

\begin{problem}
求把空间绕 $+Y$ 轴、按右手正向转过角 $\theta$ 的线性映射，并写成齐次矩阵 $R_y(\theta)$。
\end{problem}

右手正向的意思是：拇指沿 $+Y$，四指把 $+Z$ 转向 $+X$。$Y$ 坐标不变。在 $xz$ 平面上，把点写成
\[
x=r\sin\varphi,\qquad z=r\cos\varphi.
\]
$\varphi=0$ 时点在 $+Z$ 上。$\varphi$ 增加时 $x$ 从 $0$ 变为正，$z$ 从 $r$ 减小，点从 $+Z$ 转向 $+X$。转过 $\theta$ 就是把 $\varphi$ 换成 $\varphi+\theta$：
\begin{align*}
x'&=r\sin(\varphi+\theta)
=r\sin\varphi\cos\theta+r\cos\varphi\sin\theta
=x\cos\theta+z\sin\theta,\\
z'&=r\cos(\varphi+\theta)
=r\cos\varphi\cos\theta-r\sin\varphi\sin\theta
=z\cos\theta-x\sin\theta.
\end{align*}
再加上 $y'=y$，得到
\begin{equation}
R_y(\theta)=
\begin{pmatrix}
\cos\theta&0&\sin\theta&0\\
0&1&0&0\\
-\sin\theta&0&\cos\theta&0\\
0&0&0&1
\end{pmatrix}.
\label{eq:ry}
\end{equation}
检验两个基向量。$e_x=(1,0,0)$ 被送到 $(\cos\theta,0,-\sin\theta)$。小的正角度下它获得负的 $z$ 分量，也就是转向 $-Z$。$e_z=(0,0,1)$ 被送到 $(\sin\theta,0,\cos\theta)$，转向 $+X$。与右手规则一致。

程序把这张矩阵按列存成 $16$ 个数。第一列是 $(\cos\theta,0,-\sin\theta,0)$，所以数组下标 $0$ 是 $\cos\omega t$，下标 $2$ 是 $-\sin\omega t$。$\theta=\omega t$，$t$ 是秒。

\subsection{观察矩阵}

\begin{problem}
给定眼睛 $e$、目标 $c$、粗略上方向 $u_w$，且 $c-e$ 不与 $u_w$ 平行。求相机坐标，使相机位于自身原点、看向 $-Z$、右方向为 $+X$、上方向为 $+Y$。
\end{problem}

先做一组正交基。令
\begin{align*}
f&=\frac{c-e}{\lVert c-e\rVert},&
s&=\frac{f\times u_w}{\lVert f\times u_w\rVert},&
u&=s\times f.
\end{align*}
$f$ 是视线方向。$s$ 同时垂直于视线和粗略上方向，取作右方向。$u=s\times f$ 垂直于 $s$ 和 $f$，而且是单位向量，因为 $s\perp f$ 且二者都已单位化。$u$ 一般不等于 $u_w$：$u_w$ 只负责「头不要倒向哪一边」，真正写进矩阵的上方向必须和视线垂直。

世界中的点 $p$ 相对眼睛的向量是 $p-e$。把它投到这组基上。相机看向 $-Z$，所以视线方向 $f$ 在相机坐标里是 $-e_z$：
\begin{align*}
x_{\mathrm{cam}}&=s\cdot(p-e),\\
y_{\mathrm{cam}}&=u\cdot(p-e),\\
z_{\mathrm{cam}}&=(-f)\cdot(p-e)=-f\cdot p+f\cdot e.
\end{align*}
写成齐次矩阵，
\begin{equation}
V=
\begin{pmatrix}
s_x&s_y&s_z&-s\cdot e\\
u_x&u_y&u_z&-u\cdot e\\
-f_x&-f_y&-f_z&f\cdot e\\
0&0&0&1
\end{pmatrix}.
\label{eq:view}
\end{equation}

代入 $e=(1.5,1.2,2.5)$，$c=0$，$u_w=(0,1,0)$。此时 $f=-e/\lVert e\rVert$，$\lVert e\rVert=\sqrt{9.94}\approx 3.1528$。数值到四位小数是
\begin{align*}
f&\approx(-0.4758,-0.3806,-0.7930),\\
s&\approx(0.8575,0,-0.5145),\\
u&\approx(-0.1958,0.9247,-0.3264).
\end{align*}
于是
\begin{equation}
V\approx
\begin{pmatrix}
0.8575&0&-0.5145&0\\
-0.1958&0.9247&-0.3264&0\\
0.4758&0.3806&0.7930&-3.1528\\
0&0&0&1
\end{pmatrix}.
\label{eq:view-num}
\end{equation}
第四列是世界原点在相机坐标中的像。它的 $z$ 分量约等于 $-3.1528<0$，所以立方体中心在相机前方。$s\cdot e$ 在这个特殊眼睛位置上接近 $0$，第四列的 $x$ 分量因此接近 $0$，这是这组数字的巧合，不是公式的一般结论。

\subsection{透视矩阵}

\begin{problem}
相机坐标里，保留 $-f\le z\le -n$ 的对称视锥，$0<n<f$。竖直视野为 $\theta$，宽高比为 $a$。求 $4\times 4$ 矩阵 $P$，使除以 $w_c$ 之后：视锥的侧面落到 $x_n=\pm 1$、$y_n=\pm 1$，近平面落到 $z_n=-1$，远平面落到 $z_n=1$，并且 $w_c=-z$。
\end{problem}

设 $\tau=\tan(\theta/2)$。近平面 $z=-n$ 上，竖直半高是 $n\tau$，水平半宽是 $an\tau$。

要求 $w_c=-z$。点在前方时 $z<0$，所以 $w_c>0$。这决定第四行：
\[
(0,\ 0,\ -1,\ 0).
\]

侧面的条件先看 $y$。近平面右上角的方向应变成 $y_n=1$。对一般点，相似比给出
\[
y_n=\frac{y/(-z)}{\tau}=\frac{y}{\tau\,(-z)}=\frac{y/\tau}{w_c}.
\]
因此 $y_c=y/\tau$，第二行是 $(0,\ 1/\tau,\ 0,\ 0)$。同理
\[
x_n=\frac{x/(-z)}{a\tau},
\]
第一行是 $(1/(a\tau),\ 0,\ 0,\ 0)$。

深度要把 $z=-n$ 送到 $-1$，把 $z=-f$ 送到 $1$，并且 $z_c$ 只能是 $Az+B$，否则除法之前的映射不是线性的。于是
\[
z_n=\frac{Az+B}{-z}=-A-\frac{B}{z}.
\]
代入两个平面：
\begin{align*}
-A+\frac{B}{n}&=-1,\\
-A+\frac{B}{f}&=1.
\end{align*}
两式相减，
\[
B\Bigl(\frac{1}{n}-\frac{1}{f}\Bigr)=-2,
\qquad
B\Bigl(\frac{f-n}{nf}\Bigr)=-2,
\qquad
B=-\frac{2nf}{f-n}.
\]
代回第一式，
\[
A=1+\frac{B}{n}
=1-\frac{2f}{f-n}
=\frac{f-n-2f}{f-n}
=-\frac{f+n}{f-n}.
\]
第三行是 $\bigl(0,\ 0,\ -(f+n)/(f-n),\ -2nf/(f-n)\bigr)$。合起来，
\begin{equation}
P=
\begin{pmatrix}
1/(a\tau)&0&0&0\\
0&1/\tau&0&0\\
0&0&-(f+n)/(f-n)&-2nf/(f-n)\\
0&0&-1&0
\end{pmatrix}.
\label{eq:persp}
\end{equation}

代入 $\theta=50^\circ$，$a=16/9$，$n=0.1$，$f=100$。$\tau=\tan 25^\circ\approx 0.4663$，
\begin{equation}
P\approx
\begin{pmatrix}
1.2063&0&0&0\\
0&2.1445&0&0\\
0&0&-1.0020&-0.2002\\
0&0&-1&0
\end{pmatrix}.
\label{eq:persp-num}
\end{equation}

\begin{remark}
GAMES101 的推导先把视锥压成盒子，再做正交投影，近远平面取相机空间里的负 $z$。那两步乘完，再把参数改成正的距离，得到的就是 \eqref{eq:persp}。课程矩阵的第四行若是 $(0,0,1,0)$，它配套的是带符号的平面位置，不能把正的 $0.1$ 和 $100$ 直接填进去。
\end{remark}

\subsection{一个角的屏幕坐标}

\begin{problem}
$M$ 取单位矩阵时，求立方体中心和角 $q=(1/2,-1/2,1/2)$ 的标准化设备坐标。若帧缓冲是 $960\times 540$，它们的像素坐标是多少？
\end{problem}

中心 $p=0$ 乘 $V$ 得到 \eqref{eq:view-num} 的第四列。相机坐标约为 $(0,\ 0,\ -3.1528)$，齐次分量 $1$。再乘 \eqref{eq:persp-num}：
\begin{align*}
x_c&=0,\\
y_c&=0,\\
z_c&=(-1.0020)(-3.1528)+(-0.2002)\approx 2.9589,\\
w_c&=-(-3.1528)=3.1528.
\end{align*}
因此
\[
(x_n,y_n,z_n)\approx(0,\ 0,\ 0.9385).
\]
像素
\[
x_{\mathrm{pix}}=480,\qquad y_{\mathrm{pix}}=270.
\]

对角 $q$，同样先乘 $V$ 再乘 $P$，得到
\[
(x_n,y_n,z_n)\approx(0.0764,\ -0.5728,\ 0.9281).
\]
三个分量都在 $[-1,1]$ 内，$w_c>0$。像素
\[
x_{\mathrm{pix}}\approx 516.7,\qquad y_{\mathrm{pix}}\approx 115.4.
\]
它在中心的右下方。$y_{\mathrm{pix}}$ 更小，是因为视口的 $y$ 向上增加，而这个角的世界坐标 $y=-1/2$，低于中心。八个角都做过同样的检查，不旋转时和绕 $Y$ 转 $40^\circ$ 时都落在可见立方体内。程序启动时若有一个角失败，会在进入画面循环之前抛出异常。

\subsection{深度在比较什么}

\begin{problem}
两个片段落到同一像素时，保留哪一个？
\end{problem}

\eqref{eq:persp} 把 $z=-n$ 送到 $z_n=-1$，把 $z=-f$ 送到 $z_n=1$。沿视线往远处走，$z$ 变得更负，$z_n$ 增大。所以更近的面有更小的 $z_n$。默认的深度测试保留更小的那个值。清空深度时把缓冲设成 $1$，也就是先假定每一像素都被远平面占着。

背面剔除用的不是深度，而是问题 1 里的叉积。外法向朝向相机的面留下，朝内的面丢掉。深度测试和剔除回答的是两件不同的事：前者处理正面之间的遮挡，后者丢掉显然背对眼睛的三角形。

\subsection{转头：偏航与俯仰}

\begin{problem}
眼睛 $e$ 固定。哪些两个角可以表示「头朝哪」，怎样从这两个角得到视线方向，使鼠标向右是向右转、鼠标向上是抬头，并且 $\psi=0$、$\varphi=0$ 时看向 $-Z$？
\end{problem}

偏航 $\psi$ 绕世界 $+Y$ 轴。俯仰 $\varphi$ 是视线相对水平面的仰角。取
\begin{equation}
f(\psi,\varphi)=
\begin{pmatrix}
\cos\varphi\sin\psi\\
\sin\varphi\\
-\cos\varphi\cos\psi
\end{pmatrix}.
\label{eq:look}
\end{equation}
$\psi=0$、$\varphi=0$ 时，$f=(0,0,-1)$。$\varphi$ 很小且为正时，$y$ 分量为正，头抬起来。$\psi$ 很小且为正、$\varphi=0$ 时，$f\approx(\psi,0,-1)$，视线从 $-Z$ 偏向 $+X$。人站在 $+Z$ 一侧面向 $-Z$ 时，$+X$ 在右手边，所以 $\psi$ 增加就是向右转。

这和 \eqref{eq:ry} 的右手正向相反。$R_y(\psi)$ 把 $-Z$ 转到 $-X$。向右转要的是 $+X$，所以 \eqref{eq:look} 里的偏航等于把 $(0,0,-1)$ 用 $R_y(-\psi)$ 去转，再把结果抬高 $\varphi$。把俯仰写成绕相机右轴的旋转会得到同一支公式，这里直接检验分量即可。

目标点取 $e+f$。它和 $e$ 的差就是单位视线，代入 \eqref{eq:view} 得到新的 $V$。眼睛不变，变的只有 $f$。

俯仰接近 $\pm\pi/2$ 时，$f$ 与世界上方向 $(0,1,0)$ 平行，$f\times u_w=0$，\eqref{eq:view} 里给 $s$ 做归一化的那一次除法没有定义。程序把 $\varphi$ 夹在 $\pm 89^\circ$。

启动时要求 $f$ 指向原点，也就是与 $-e$ 同向。此时 $\sin\varphi=f_y$，并且 $\cos\varphi>0$，故
\[
\varphi=\arcsin(f_y),\qquad
\psi=\operatorname{atan2}(f_x,-f_z).
\]
后一式是因为 $f_x=\cos\varphi\sin\psi$，$-f_z=\cos\varphi\cos\psi$，正的 $\cos\varphi$ 不改变 $\operatorname{atan2}$ 的两个自变量之比。对 $e=(1.5,1.2,2.5)$，
\[
f=\frac{-e}{\lVert e\rVert},\qquad
\psi\approx -31.0^\circ,\qquad
\varphi\approx -22.4^\circ.
\]
用这两个角重建的 $V$ 必须和直接看向原点的 $V$ 逐项一致，容差 $10^{-4}$。不一致就说明 \eqref{eq:look} 的符号和观察矩阵不是同一套约定。

鼠标给出的是像素增量，不是角度。灵敏度 $\sigma=0.002$ 弧度每像素。GLFW 的光标 $x$ 向右增加，$y$ 向下增加。因此
\[
\psi\leftarrow\psi+\sigma\,\Delta x,\qquad
\varphi\leftarrow\varphi-\sigma\,\Delta y.
\]
第一帧的 $\Delta x,\Delta y$ 丢掉。光标刚被隐藏时，系统常会送来一次虚假的大跳跃。若驱动允许，使用原始鼠标运动，避开操作系统的指针加速。

\subsection{在水平面上走}

\begin{problem}
眼睛可以移动之后，W 应该沿着哪一个方向走？抬头时按下 W，人该不该离开水平面？同时按住 W 和 D，速率该不该变成单键的 $\sqrt{2}$ 倍？
\end{problem}

俯仰已经放在视线 $f(\psi,\varphi)$ 里。若位移直接沿 $f$，抬头再按 W 就会沿视线飞起来。这一步还没有重力，也没有脚底的平面；先规定人站在启动时的高度上，只在 $y=\text{常数}$ 的平面里走。水平前方向是 \eqref{eq:look} 丢掉俯仰之后的单位向量
\begin{equation}
h(\psi)=\begin{pmatrix}\sin\psi\\ 0\\ -\cos\psi\end{pmatrix}.
\label{eq:walk-forward}
\end{equation}
$\psi=0$ 时 $h=(0,0,-1)$。水平右方向取与 $h$ 垂直、且 $\psi=0$ 时指向 $+X$ 的单位向量
\begin{equation}
r(\psi)=\begin{pmatrix}\cos\psi\\ 0\\ \sin\psi\end{pmatrix}.
\label{eq:walk-right}
\end{equation}
检验 $h\cdot r=\sin\psi\cos\psi-\cos\psi\sin\psi=0$，并且 $\lVert h\rVert=\lVert r\rVert=1$。$\psi=0$ 时 $r=(1,0,0)$，正是面向 $-Z$ 时的右手边。

令 $w\in\{-1,0,1\}$ 表示 S、空、W，$d\in\{-1,0,1\}$ 表示 A、空、D。愿望方向
\[
u=w\,h+d\,r.
\]
只按 W 时 $u=h$，长度是 $1$。同时按 W 和 D 时 $u=h+r$，长度是 $\sqrt{2}$。若直接用 $u$ 乘速率，斜着走会更快。长度非零时先归一化。速率 $s=2$（每秒世界单位），帧间隔为 $\Delta t$ 秒，且把 $\Delta t$ 截断到不超过 $0.05$，避免一次卡顿把人送出很远。眼睛更新为
\begin{equation}
e\leftarrow e+s\Delta t\,\frac{u}{\lVert u\rVert},\qquad \lVert u\rVert>0.
\label{eq:walk}
\end{equation}
$y$ 分量始终是 $0$，所以 $e_y$ 不变。

$\psi=0$、只按住 W 一秒钟、且 $\Delta t$ 没有被截断时，位移是 $s h=(0,0,-2)$。眼睛从 $(1.5,1.2,2.5)$ 走到 $(1.5,1.2,0.5)$，高度仍是 $1.2$。立方体占据 $|x|,|y|,|z|\le 1/2$。这条路线的 $x=1.5$ 始终在立方体外，所以这一秒还撞不上它。若先把朝向转到正对立方体再前进，眼睛的水平位置会进入 $|x|,|z|\le 1/2$。眼睛的高度 $1.2$ 仍在立方体上方。下一节用身体的盒子，而不是眼睛这个点，来决定挡不挡。

\subsection{身体和立方体}

\begin{problem}
眼睛固定在 $y=1.2$，立方体只占据 $|y|\le 1/2$。为什么用眼睛这个点做碰撞，人可以在立方体正上方走过去却永远不会被挡住？身体该取多大，才既能在出发时不重叠，又能在走到立方体上方时被顶住？
\end{problem}

立方体是中心在原点、半边长 $1/2$ 的轴对齐盒子。眼睛 $(0,1.2,0)$ 满足 $|1.2|>1/2$，所以这个点在立方体外。只检测点是否在盒子内，水平移动永远碰不到它。

身体也取轴对齐盒子，中心放在眼睛上。半宽（也是半深）$a=0.3$，半高 $b=0.9$。身体在竖直方向占据
\[
[1.2-0.9,\ 1.2+0.9]=[0.3,\ 2.1].
\]
立方体占据 $[-1/2,1/2]$。两段的交集是 $[0.3,1/2]$，不是空集。出发位置 $(1.5,1.2,2.5)$ 在水平方向离原点更远，后面会算出它不重叠。

两个轴对齐盒子重叠，当且仅当每个轴上的区间都重叠。中心差的绝对值小于半边长之和：
\begin{equation}
|x|<a+\tfrac12,\qquad |y|<b+\tfrac12,\qquad |z|<a+\tfrac12.
\label{eq:aabb}
\end{equation}
代入数字，$|x|<0.8$，$|y|<1.4$，$|z|<0.8$。当前 $y=1.2<1.4$，竖直条件始终成立，碰撞退化为水平面上的正方形 $|x|<0.8$、$|z|<0.8$。等号不算重叠，所以身体刚好贴在这个正方形边界上时允许停住。

先把 $x$ 方向的位移加上，若 \eqref{eq:aabb} 成立就把 $x$ 退回；再对 $z$ 做同样的事。斜着撞上一个角时，被挡住的那一轴取消，另一轴还在，人会沿着棱滑过去。两轴一起加再整体退回，会把沿着墙的那部分速度也吃掉。

自检四个命题，有一个失败就不进入画面：眼睛 $(0,1.2,0)$ 不在立方体内部；以这个眼睛为中心的身体与立方体重叠；出发位置不重叠；$x=0.8$、$z=0$ 的贴边位置不重叠。

\subsection{地板、重力和方块编号}

\begin{problem}
只有原点一块实心方块时，人从 $(1.5,1.2,2.5)$ 松开按键会一直下落。怎样用整数编号表示许多方块，使人能落在一块地板上，并且编号 $(0,0,0)$ 仍然就是原来那一块？
\end{problem}

方块 $(i,j,k)$ 占据
\[
[i-\tfrac12,i+\tfrac12]\times[j-\tfrac12,j+\tfrac12]\times[k-\tfrac12,k+\tfrac12].
\]
原来的立方体就是 $(0,0,0)$。它的碰撞公式 \eqref{eq:aabb} 是这个定义在 $i=j=k=0$ 时的特例。身体与方块 $(i,j,k)$ 重叠当且仅当
\[
|x-i|<0.8,\qquad |y-j|<1.4,\qquad |z-k|<0.8.
\]
数组只保存 $i,j,k\in[-4,4]$。这个范围之外的查询一律当空气。地板是 $j=-1$、$i,k\in[-2,2]$ 的 $25$ 块，再加上 $(0,0,0)$，共 $26$ 块实心。

竖直速度 $v$ 每帧先被重力改变，再用来移动眼睛：
\[
v\leftarrow v-20\Delta t,\qquad e_y\leftarrow e_y+v\Delta t.
\]
若移动后身体与任何实心方块重叠，就把 $e_y$ 退回，并将 $v$ 设为 $0$。水平方向仍按轴分开处理。

地板顶面在 $y=-1/2$。身体半高 $0.9$，贴住顶面时眼睛高度是
\[
-1/2+0.9=0.4.
\]
此时 $y$ 方向刚好贴上，不算重叠。眼睛再低到 $0.39$，身体底端进入地板，重叠成立。出发高度 $1.2$ 在地板上方，所以开场会落下，大约四分之一秒后停在 $y=0.4$。$(0,0,0)$ 那一块高出地板，顶面在 $y=1/2$，人从侧面走过去会被它挡住，这一章还不能跳上去。

模型矩阵不再随时间旋转。方块 $(i,j,k)$ 的模型矩阵是把立方体平移到 $(i,j,k)$ 的
\[
T=
\begin{pmatrix}
1&0&0&i\\
0&1&0&j\\
0&0&1&k\\
0&0&0&1
\end{pmatrix}.
\]
若视觉上继续旋转而碰撞盒不动，人会看见自己穿过方块的角，或被一块已经转开的面挡住。

\subsection{挖掉看着的那一块}

\begin{problem}
左键按下时，该把哪一块实心方块改成空气？为什么不能把眼睛所在的格子四舍五入之后就挖掉？
\end{problem}

视线是已经单位化的方向 \(d\)。射线上的点是
\[
p(t)=e+td,\qquad 0\le t\le 5.
\]
\(t\) 就是沿视线走了多远，最远 \(5\)。方块 \((i,j,k)\) 在世界里占 \([i-1/2,i+1/2]\)。把它平移到网格坐标 \(g=p+(1/2,1/2,1/2)\)，这块方块就变成半开区间 \([i,i+1)\)。眼睛一开始落在
\[
i_0=\lfloor e_x+1/2\rfloor,
\]
\(j_0,k_0\) 同样计算。这个格子通常是空气。

沿每个轴问下一次碰到格子边界要走多少 \(t\)。\(d_x>0\) 时，\(x\) 方向的下一道边界是 \(i_0+1\)，
\[
t_x=\frac{(i_0+1)-g_x}{d_x}.
\]
\(d_x<0\) 时边界是 \(i_0\)，分母为负，\(t_x\) 仍为正。\(d_x=0\) 时这一轴永远不跨格，\(t_x\) 记成无穷。\(y,z\) 同样。每跨过一格，该轴的 \(t\) 再增加 \(1/|d|\)，因为格子边长是 \(1\)。

每一轮取 \(t_x,t_y,t_z\) 里最小的那个，沿对应的轴走进相邻格子。若两个相等，说明射线正好穿过棱或角，这些轴一起前进，避免漏掉斜着贴住的方块。走进新格子后若它是实心的，并且 \(t\le 5\)，这就是要挖的块。把数组里的那一格写成空气。画面和碰撞都只问这个数组，所以块会消失，人也能走过去。

一次按下只挖一块。按住不放时，后面的帧不再重复触发。否则一秒钟会沿视线挖空一串。

可用的检验：从 \((0,0.4,2)\) 朝 \(-Z\) 看，方向 \((0,0,-1)\)。第一块实心是 \((0,0,0)\)，碰上它的面 \(z=1/2\) 时 \(t=2-1/2=1.5\)。同一点朝下看，方向 \((0,-1,0)\)，第一块是地板 \((0,-1,2)\)。朝上看，\(5\) 的距离内没有实心块。

\subsection{关掉再打开}

\begin{problem}
挖掉的方块只存在内存里的布尔数组中。进程一结束，数组就没了。怎样把「哪些格子是实心的」写成文件，使下一次启动读回来的是挖过之后的世界，而不是重新生成的那 \(26\) 块？
\end{problem}

文件第一行固定为
\[
\texttt{cubic-dawn-blocks 1}
\]
后面每一行是三个整数 \(i\ j\ k\)，表示这一格为实心。没有出现的格子是空气。范围仍是 \([-4,4]\)。读入时先把数组全部清成空气，再按行设成实心。所以文件是世界的完整名单，不是「在默认世界上打过的补丁」。若只把挖掉的块追加成某种删除记录，又忘了重写剩余的块，旧块会在下次启动时回来。

写出的顺序按 \(j\)、\(k\)、\(i\) 从小到大，同一份世界每次得到同一个文件。挖掉一块之后立刻重写整份文件，退出时再写一次。人的坐标和朝向不在这个文件里，下次仍从 \((1.5,1.2,2.5)\) 出生，然后再落到地板上。

没有这个文件时，先按下一节的规则生成地形，自检通过后再写出文件。已有同一版本的文件时，自检仍先在新生成的地形上做完，再整份读入。旧的 \(26\) 块手摆存档版本号不同，会被新地形覆盖。

\subsection{种子算出来的地面}

\begin{problem}
手摆的 \(5\times 5\) 地板和单独一块柱子，仍然是一份图形学作业里的场景：格子、射线、存档都有了，世界却是人写死的。怎样让同一颗种子每次长出同一片高低起伏的地面？
\end{problem}

这一节的数学比前面的矩阵长，而且每一步都改得了最终高度。下面按计算顺序写，不把散列、平滑和插值并成一句。

\paragraph{随机数不能用。}
\texttt{Math.random} 每次启动给一套新数。挖过的存档对不上下一回长出来的山。这里要的是函数 \(r(p,q)\)：同一颗种子、同一对整数，永远得到同一个小数。

\paragraph{32 位散列，用 \((p,q)=(0,0)\) 走一遍。}
种子 \(s=20261116\)。四个乘数都是奇数，没有几何意义，只负责把三个整数搅进 \(32\) 位的每一位。运算一律留在 \(32\) 位补码里，也就是随时对 \(2^{32}\) 取余。逻辑右移记作 \(\gg\)，高位补 \(0\)，不是把负数的符号位继续补 \(1\)。
\[
n=374761393\,p+668265263\,q+1442695041\,s.
\]
\((0,0)\) 只剩下种子那一项：
\[
n_0=20261116\times 1442695041=464620284.
\]
这个数小于 \(2^{32}\)，没有发生回绕。搅动分两下。先与自己右移 \(13\) 位的结果异或，再乘 \(1274126177\)：
\begin{align*}
n_1&=464620284\oplus (464620284\gg 13)=464607088,\\
n_2&=464607088\times 1274126177\bmod 2^{32}=3287331184.
\end{align*}
再异或一次右移 \(16\) 位：
\[
n_3=3287331184\oplus (3287331184\gg 16)=3287315072.
\]
最低 \(8\) 位是 \(3287315072\bmod 256=128\)。格子上的高度小数是
\[
r(0,0)=\frac{128}{255}\approx 0.501961.
\]
分母是 \(255\)，不是 \(256\)。字节 \(255\) 会得到 \(r=1\)。于是 \(\nu\) 可以等于 \(1\)，下面的 \(\lfloor 4\nu\rfloor\) 可以等于 \(4\)。这一种子在我们抽查的点上没有碰上 \(1\)，但公式允许第五级高度。只留 \(8\) 位，是因为地形只需要一个 \([0,1]\) 里的小数；另外 \(24\) 位丢掉，不会让同一对 \((p,q)\) 变出第二个数。

\paragraph{补码：很长的那个数为什么能是负数。}
\(3287331184\) 看起来是一个普通的正整数。程序里的 \texttt{int} 只有 \(32\) 位，没有单独的符号位存放处。补码的规则是：最高位是 \(0\)，按普通二进制读，范围是 \(0\) 到 \(2^{31}-1=2147483647\)；最高位是 \(1\)，同一个位模式表示
\[
\text{位模式}-2^{32}.
\]
先看 \(4\) 位，一共只能写下 \(0000\) 到 \(1111\)。后一半的最高位是 \(1\)：
\begin{align*}
0000&=0,&
0111&=7,\\
1000&=-8,&
1011&=-5,\\
1111&=-1.
\end{align*}
检验 \(1011\)。位模式是 \(11\)，\(11-2^4=11-16=-5\)。把 \(1011\) 按位取反得 \(0100\)，再加 \(1\) 得 \(0101=5\)，所以它是 \(-5\)。这就是「取反加一」：负数的补码，取反加一之后回到它的绝对值。

\(32\) 位是同一件事，只是模数换成 \(2^{32}=4294967296\)。\(3287331184>2147483647\)，最高位是 \(1\)，写成十六进制是 \texttt{0xC3F0B170}。Java 的 \texttt{int} 把它读成
\[
3287331184-4294967296=-1007636112.
\]
后面的 \(n_3=3287315072\) 同样是负数，等于 \(-1007652224\)。散列公式并不在乎这个读法。异或、逻辑右移、乘法都作用在位上。最低 \(8\) 位用的是按位与 \(255\)，也就是留下 \texttt{xxxxxxxx} 的最后八位。\(n_3\) 的最后八位仍是 \(128\)。若改成对负数做 Java 的 \texttt{\%} 运算，余数可以是负的，\(r(p,q)\) 就会掉出 \([0,1]\)。位运算不受补码读法的影响，取余不一定。

\paragraph{为什么不直接把 \(r(p,q)\) 当成这一列的高度。}
若每一列各抽一个独立小数，相邻列之间没有关系，地面是噪声点，不是山坡。所以随机数只放在边长为 \(4\) 的粗格子角上，格子内部用插值填。方块 \((x,z)\) 所在的格子是
\[
p=\Bigl\lfloor\frac{x}{4}\Bigr\rfloor,\qquad
f_x=\frac{x}{4}-p,
\]
\(q,f_z\) 同样。\(f_x,f_z\in[0,1)\)。\((0,0)\) 正好落在角上，\(f_x=f_z=0\)，插值结果就是 \(r(0,0)\)，不用看另外三角。

\paragraph{平滑权重，关键在导数。}
令
\[
u(t)=t^2(3-2t)=3t^2-2t^3.
\]
它仍然把 \([0,1]\) 映到 \([0,1]\)，而且
\[
u(0)=0,\qquad u(1)=1,\qquad u\bigl(\tfrac12\bigr)=\tfrac12.
\]
线性权重 \(t\) 也满足这三式。差别在导数：
\[
u'(t)=6t-6t^2=6t(1-t),\qquad u'(0)=u'(1)=0.
\]
线性权重的导数恒为 \(1\)。一块格子的右边界和下一块的左边界接在同一组角点上，高度连续；但线性权重在边界上的斜率是 \(1\)，山坡会折出一条棱。\(u'\) 在两端为 \(0\)，穿过粗格子的整数边时，坡度先回到 \(0\) 再离开。靠近边缘时它也比线性更贴着起点：\(u(0.1)=0.028\)，而线性权重已经是 \(0.1\)。

\paragraph{\((2,3)\) 的四角，一步不省。}
\[
\frac{2}{4}=0.5,\qquad \frac{3}{4}=0.75,
\]
所以 \(p=0\)，\(q=0\)，\(f_x=1/2\)，\(f_z=3/4\)。
\[
u(f_x)=u\bigl(\tfrac12\bigr)=\tfrac12,
\]
\[
u(f_z)=u\bigl(\tfrac34\bigr)
=\Bigl(\tfrac34\Bigr)^2\Bigl(3-2\cdot\tfrac34\Bigr)
=\tfrac{9}{16}\cdot\tfrac{3}{2}
=\tfrac{27}{32}=0.84375.
\]
四个角用同一套散列。最低 \(8\) 位和对应的小数是
\begin{align*}
r(0,0)&=\frac{128}{255}\approx 0.501961,\\
r(1,0)&=\frac{68}{255}\approx 0.266667,\\
r(0,1)&=\frac{177}{255}\approx 0.694118,\\
r(1,1)&=\frac{74}{255}\approx 0.290196.
\end{align*}
先沿 \(x\) 插值。\(u(f_x)=1/2\)，所以就是两角的算术平均：
\begin{align*}
a&=\frac{1}{2}r(0,0)+\frac{1}{2}r(1,0)
=\frac{128+68}{2\times 255}
=\frac{98}{255}\approx 0.384314,\\[4pt]
b&=\frac{1}{2}r(0,1)+\frac{1}{2}r(1,1)
=\frac{177+74}{2\times 255}
=\frac{251}{510}\approx 0.492157.
\end{align*}
再沿 \(z\)：
\begin{align*}
\nu
&=a\bigl(1-u(f_z)\bigr)+b\,u(f_z)\\
&=0.384314\times 0.15625+0.492157\times 0.84375\\
&\approx 0.475306.
\end{align*}
顶面高度只保留整数级：
\[
h(2,3)=\lfloor 4\nu\rfloor=\lfloor 1.901\rfloor=1.
\]
\(4\nu\) 的小数部分 \(0.901\) 被丢掉。地形只有五档可能的整数，中间的坡度靠相邻列的 \(h\) 不同来形成，不靠这一列内部的小数。

\paragraph{这一列到底放哪些方块。}
\(h=1\) 是最上面一块的编号。厚度固定为 \(4\)，所以实心的是
\[
y=1,0,-1,-2.
\]
再往上的 \(y=2\) 是空气，再往下的 \(y=-3\) 也是空气。人站在这块中心时，脚底贴着顶面 \(y=1+1/2=1.5\)。眼睛在身体中心，半高 \(0.9\)，贴住时
\[
e_y=1.5+0.9=2.4.
\]
程序再加 \(0.25\)，从 \(2.65\) 落下去。岛的范围是 \(x,z\in[-10,10]\)，共 \(21\times 21\) 列，每列 \(4\) 块，合计 \(1764\)。数组能存到 \([-16,16]\)，岛外是空气。

同一颗种子的另外两个抽查：\(h(0,0)=2\)，因为这一点就在角上，\(\nu=r(0,0)\)，\(\lfloor 4\times 0.501961\rfloor=\lfloor 2.007\rfloor=2\)。\(h(10,10)=0\)。

\paragraph{负坐标。}
\(-10/4=-2.5\)。向下取整
\[
\lfloor -2.5\rfloor=-3,\qquad f=(-2.5)-(-3)=0.5.
\]
向零取整得到 \(-2\)，小数部分会变成 \(-0.5\)，不再落在 \([0,1)\)。格子编号差 \(1\) 之后，岛的西边和南边接到另一套角点上，\(x=-8\) 或 \(z=-8\) 附近会裂开。\((-10,-10)\) 用的角是 \((-3,-3)\) 及其加一的三个邻居，不是 \((-2,-2)\)。

\subsection{方块有了编号}

\begin{problem}
\(1764\) 块现在只有「在」或「不在」。挖掉哪一块、脚被哪一块挡住，都不区分草、泥土和石头。编号和属性各管什么？
\end{problem}

空气的编号是 \(0\)，草是 \(1\)，泥土是 \(2\)，石头是 \(3\)。数组里存的是这个编号，不再是一个布尔值。属性表只有两列：是否挡人，以及乘到顶点颜色上的三个小数。

\begin{center}
\begin{tabular}{c|c|c}
编号 & 挡住人 & 颜色 \((r,g,b)\) \\
\hline
\(0\) 空气 & 否 & 不画 \\
\(1\) 草 & 是 & \((0.36,0.70,0.28)\) \\
\(2\) 泥土 & 是 & \((0.58,0.38,0.22)\) \\
\(3\) 石头 & 是 & \((0.62,0.64,0.66)\)
\end{tabular}
\end{center}

顶点颜色仍然区分六个面。片元颜色是顶点颜色和这张表里的颜色逐分量相乘，所以同一块的六个面还有明暗差别，不同种类的整体色不同。这一步没有贴图。贴图是把表面上的小数坐标再映射到一张图，属性表里现在还没有那一列。

一列的顶面用草，顶面下面一块用泥土，再往下用石头。\((2,3)\) 的顶面 \(h=1\)，所以 \(y=1\) 是草，\(y=0\) 是泥土，\(y=-1\) 和 \(y=-2\) 是石头。存档每一行变成四个整数 \(i\ j\ k\ \text{编号}\)，版本头是 \texttt{cubic-dawn-blocks 3}。旧文件没有编号，读不进来，会按这个规则重新生成。

\subsection{空格起跳}

\begin{problem}
脚已经能落在方块上，人却迈不上高一格的坡。怎样用一次向上的初速度让身体越过一格，同时避免在空中连按空格就飞起来？
\end{problem}

竖直方向已经有
\[
v\leftarrow v-20\Delta t.
\]
起跳不增加一项新的力，只在离地的那一帧把 \(v\) 设成 \(7.1\)，这一帧不再先减重力。下一帧人已经离开地面，重力照常作用。到达最高点时 \(v=0\)，上升高度是
\[
h=\frac{v_0^2}{2g}=\frac{7.1^2}{2\times 20}\approx 1.26.
\]
一格方块的高度是 \(1\)，所以这次起跳过得去，也过不了两格叠在一起的墙。

「在地上」不能看 \(v\) 是不是 \(0\)。落地那一帧把重叠的位移退回后，人可能还悬在表面上方一点点。判断是把眼睛再降低 \(0.05\)，身体若因此碰到实心方块，才允许起跳。空中这个探测碰不到方块，空格没有效果。按住不放时，只在按下的第一帧给速度；落地后若键已经松开再按下，才能跳下一次。

\begin{insight}
这一节里最容易被当成装饰的是 \(u(t)=3t^2-2t^3\)。它在 \(0\)、\(1/2\)、\(1\) 三点上和线性权重 \(t\) 一致，所以只检查这三个点看不出差别。要看的是 \(u'(0)=u'(1)=0\)。没有这个条件，粗格子的每条边上都会留下一条折棱。
\end{insight}

\section{关键 Insight}

\begin{insight}
齐次坐标把两种失败收成同一种成功。平移失败，是因为线性映射留不住一个加项；透视失败，是因为近大远小不是线性的。多出来的 $w$ 一次只解决一件事：平移时用它等于 $1$ 去承接第四列，透视时用它等于 $-z$ 去承接随后的除法。管线看起来是三张矩阵相乘，几何上却是「线性部分排队，非线性只留下一次除法」。
\end{insight}

\section{验证题与错因分析}

\begin{exercise}[旋转把 $+X$ 送到哪里]
取 $\theta=\pi/2$。用 \eqref{eq:ry} 计算 $R_y(\theta)$ 作用在 $(1,0,0)$ 上的结果，并说明它相对于 $+Z$ 的位置。
\end{exercise}

\begin{exercise}[视线轴上的点]
不计算 $s$ 和 $u$ 的具体分量，说明世界原点的 $x_n$ 和 $y_n$ 为什么一定是 $0$。再说明：若把目标改成 $(0,1,0)$，这个结论还在不在。
\end{exercise}

\begin{exercise}[向右看]
$\varphi=0$。用 \eqref{eq:look} 计算 $\psi=\pi/2$ 时的视线。鼠标向右移动 $100$ 像素、灵敏度为 $0.002$ 时，$\psi$ 增加多少弧度？
\end{exercise}

\begin{exercise}[走一秒]
$\psi=0$，只按住 W，速率为 $2$，一秒钟内帧间隔都没有被截断。眼睛的位移是多少？同时按住 W 和 D 时，若不做 \eqref{eq:walk} 里的归一化，这一秒的位移长度是多少？
\end{exercise}

\begin{exercise}[眼睛和身体]
眼睛放在 $(0,1.2,0)$。它在不在立方体内部？以它为中心、半宽 $0.3$、半高 $0.9$ 的身体与立方体重叠吗？$x=0.8$、$y=1.2$、$z=0$ 呢？
\end{exercise}

\begin{exercise}[落在地板上]
地板方块 $j=-1$ 的顶面在 $y=-1/2$，身体半高是 $0.9$。贴住时眼睛的 $y$ 是多少？$y$ 再减小 $0.01$ 时，重叠判断为什么变成真？
\end{exercise}

\begin{exercise}[朝柱子看]
从 \((0,0.4,2)\) 沿 \((0,0,-1)\) 看，最远距离 \(5\)。第一块实心方块的编号是什么？射线走到它表面时 \(t\) 是多少？
\end{exercise}

\begin{exercise}[只剩一行]
文件第一行是版本头，第二行只有 \(0\ -1\ 0\)。读入之后，\((0,0,0)\) 和 \((0,-1,0)\) 分别是实心还是空气？默认的另外 \(24\) 块地板还在吗？
\end{exercise}

\begin{exercise}[平滑权重]
计算 \(u(0)\)、\(u(1)\)、\(u(1/2)\) 和 \(u'(0)\)、\(u'(1)\)。种子 \(20261116\) 下，已知
\[
r(0,0)=\frac{128}{255},\;
r(1,0)=\frac{68}{255},\;
r(0,1)=\frac{177}{255},\;
r(1,1)=\frac{74}{255},
\]
且 \((2,3)\) 的分数坐标是 \(f_x=1/2\)、\(f_z=3/4\)。求 \(a\)、\(b\)、\(\nu\) 和顶面高度。最下面一块的 \(y\) 是多少？
\end{exercise}

\begin{exercise}[补码和草皮]
\(4\) 位模式 \(1011\) 表示的整数是多少？\(32\) 位模式 \(3287331184\) 作为 Java 的 \texttt{int} 是正还是负，数值是多少？列 \((2,3)\) 的 \(y=1\)、\(y=0\)、\(y=-2\) 三个方块编号分别是什么？
\end{exercise}

\begin{exercise}[起跳高度]
初速度 \(7.1\)，重力加速度 \(20\)。最高点比起跳点高多少？这个高度够不够越过一格方块？人在空中时，为什么再按空格不会把速度再设成 \(7.1\)？
\end{exercise}

\begin{exercise}[近平面上的半高]
$\theta=50^\circ$，$n=0.1$。求近平面与视锥相交的矩形的竖直半高。一个相机坐标为 $(0,\ 0.2,\ -0.1)$ 的点，是否还在视锥内部？
\end{exercise}

\textbf{参考答案}：

练习 1。$\cos(\pi/2)=0$，$\sin(\pi/2)=1$，所以
\[
R_y(\pi/2)\begin{pmatrix}1\\0\\0\\1\end{pmatrix}
=\begin{pmatrix}0\\0\\-1\\1\end{pmatrix}.
\]
$+X$ 上的点被送到 $-Z$。这与推导一致：正向旋转先把 $+Z$ 送向 $+X$，转过直角后，$+X$ 就占据原来 $-Z$ 的方向。

练习 2。目标是原点时，视线 $f$ 与 $-e$ 平行，原点相对眼睛的向量 $-e$ 也与 $f$ 平行。它在 $s$ 和 $u$ 上的投影都是 $0$，因为这两组向量都垂直于 $f$。故 $x_{\mathrm{cam}}=y_{\mathrm{cam}}=0$。$P$ 的前两行不把 $z$ 混进 $x_c,y_c$，除法之后 $x_n=y_n=0$。

目标改成 $(0,1,0)$ 之后，原点不再在视线上。$V$ 的第四列不再只有 $z$ 分量，$x_n$ 或 $y_n$ 一般不为 $0$。立方体中心会离开画面中心。

练习 3。$\varphi=0$、$\psi=\pi/2$ 时，
\[
f=(1,0,0).
\]
视线指向 $+X$，也就是起步时面向 $-Z$ 的人的右手边。$100$ 像素使
\[
\Delta\psi=0.002\times 100=0.2
\]
弧度，约 $11.5^\circ$。

练习 4。由 \eqref{eq:walk-forward}，$\psi=0$ 时 $h=(0,0,-1)$。位移是
\[
2\cdot h=(0,0,-2).
\]
眼睛的 $y$ 不变。W 与 D 同时按下时 $u=h+r$，长度为 $\sqrt{2}$。不归一化的位移长度是 $2\sqrt{2}\approx 2.828$，归一化之后仍是 $2$。

练习 5。$|1.2|>1/2$，眼睛在立方体外。身体的竖直区间 $[0.3,2.1]$ 与 $[-1/2,1/2]$ 相交，且 $|x|=|z|=0<0.8$，所以重叠。$x=0.8$ 不满足严格小于 $0.8$，贴边不算重叠。

练习 6。贴住时
\[
e_y=-1/2+0.9=0.4.
\]
身体底端是 $0.4-0.9=-0.5$，与地板顶面相等。重叠用的是严格小于，所以贴住为假。再降 $0.01$ 后底端是 $-0.51$，进入方块内部，三个轴的不等式同时成立。

练习 7。格子 \(z=0\) 的靠近眼睛的一面在世界坐标 \(z=1/2\)。
\[
t=2-\tfrac12=1.5.
\]
进入的格子是 \((0,0,0)\)，它是实心的。途中 \(z=2\) 和 \(z=1\) 的 \(y=0\) 都是空气。

练习 8。读入会先清空数组。文件里只有 \((0,-1,0)\)，所以它是实心，\((0,0,0)\) 是空气。另外 \(24\) 块没有写进行，不会从默认世界里继承过来。

练习 9。\(u(0)=0\)，\(u(1)=1\)，\(u(1/2)=1/2\)。\(u'(t)=6t(1-t)\)，所以 \(u'(0)=u'(1)=0\)。
\[
u(3/4)=27/32=0.84375.
\]
\(f_x=1/2\) 时 \(a\)、\(b\) 是两对角的平均：
\[
a=\frac{98}{255}\approx 0.384314,\qquad
b=\frac{251}{510}\approx 0.492157.
\]
\[
\nu\approx 0.384314\times 0.15625+0.492157\times 0.84375\approx 0.475306,
\]
\[
h=\lfloor 4\times 0.475306\rfloor=\lfloor 1.901\rfloor=1.
\]
最下面一块是 \(1-3=-2\)。

练习 10。\(1011\) 的位模式是 \(11\)，\(11-16=-5\)。\(3287331184>2^{31}-1\)，所以是负数：
\[
3287331184-2^{32}=-1007636112.
\]
列 \((2,3)\) 的顶面是 \(y=1\)，编号 \(1\)（草）；\(y=0\) 是编号 \(2\)（泥土）；\(y=-2\) 是编号 \(3\)（石头）。

练习 11。
\[
h=\frac{7.1^2}{40}=1.26025.
\]
大于 \(1\)，够越过一格，不够越过两格。空中把眼睛降低 \(0.05\) 仍碰不到方块，起跳条件不成立。

练习 12。竖直半高是
\[
n\tau=0.1\tan 25^\circ\approx 0.04663.
\]
点 $(0,0.2,-0.1)$ 就在近平面上，但 $|y|=0.2>0.04663$，已经超出视锥的上侧面。它的 $y_n$ 会大于 $1$，被裁掉。近，不等于看得见。

\textbf{典型错因分析}：

\begin{pitfall}[错因 1：用正的 $z$ 去除]
把 $x_n$ 写成 $x/z$，并且相机前方仍取 $z<0$。此时分母为负，画面左右颠倒，再与 $w_c=-z$ 的矩阵混用时，$w_c$ 的符号和除法的符号会有一边重复、一边缺失。推导里 $w_c=-z$ 与 $x/(-z)$ 是同一次选择，必须成对出现。
\end{pitfall}

\begin{pitfall}[错因 2：近远平面的符号用错层]
课程推导中的平面位置是负的相机坐标 $-0.1$ 和 $-100$。程序接口里的 $n,f$ 是正的距离。把正数代进第四行为 $(0,0,1,0)$ 的课程矩阵，近平面的像不再是 $z_n=-1$，立方体通常整颗消失。应使用 \eqref{eq:persp}，或先把正距离写成负的平面位置再进入课程公式。
\end{pitfall}

\begin{pitfall}[错因 3：用八个角携带面颜色]
八个数的是位置。颜色属于面。一个角属于三个面，就有三种颜色。只存八个位置时，三个面不得不共用一个颜色，转起来不能把面分开。$36$ 数的是带属性的三角形顶点，不是角。
\end{pitfall}

\begin{pitfall}[错因 4：鼠标的 $y$ 和抬头同号]
光标向下时 $\Delta y>0$。若把俯仰写成 $\varphi\leftarrow\varphi+\sigma\Delta y$，鼠标上移会低头。画面坐标的 $y$ 向下，世界的 $+Y$ 向上，这两个 $y$ 不是同一个数。
\end{pitfall}

\begin{pitfall}[错因 5：用带俯仰的视线走路]
$f$ 的 $y$ 分量是 $\sin\varphi$。沿 $f$ 做位移时，抬头按 W 会离开水平面。这一步的前进方向是 \eqref{eq:walk-forward}，俯仰只影响看见什么，不影响脚往哪迈。
\end{pitfall}

\begin{pitfall}[错因 6：用眼睛那个点去撞立方体]
眼睛高度 $1.2$ 在立方体顶面 $0.5$ 之上。点在盒子内的检测永远是假，人会从立方体上方穿行。要检测的是身体区间和立方体区间在三个轴上是否同时重叠。
\end{pitfall}

\begin{pitfall}[错因 7：画面在转，碰撞盒不转]
模型矩阵若仍是 $R_y(\omega t)$，方块的角在画面上移动，碰撞用的盒子却留在原地。人要么穿过看得见的角，要么被一块已经转开的空气挡住。世界里的方块用平移矩阵 $T$，碰撞盒的中心就是 $(i,j,k)$。
\end{pitfall}

\begin{pitfall}[错因 8：按住不放就连挖]
左键的状态若每帧都处理，按住的一秒会沿视线挖掉许多块。只在「这一帧按下、上一帧没按下」时挖一块。
\end{pitfall}

\begin{pitfall}[错因 9：把删除记成补丁]
若文件只追加「挖掉了哪一块」，读入时又先生成默认的 \(26\) 块再应用补丁，格式必须同时记得补丁是否已经用过。现在的文件是剩余实心块的完整名单。少写一行，那一块就是空气。
\end{pitfall}

\begin{pitfall}[错因 10：负坐标用向零取整]
\(-10/4=-2.5\)。向零取整得到 \(-2\)，向下取整得到 \(-3\)。格子编号差 \(1\)，岛的一侧会接到另一套随机数上，高度在 \(x=-8\) 或 \(z=-8\) 附近出现一条裂缝。
\end{pitfall}

\section{进一步思考}

\begin{enumerate}
\item 若模型矩阵里出现各轴倍数不同的缩放，法向还是用 $R$ 去乘吗？找出一个二维例子：把单位圆沿 $x$ 拉长为椭圆，原来的法向 $(0,1)$ 被同样的缩放送到哪里，真正的法向又在哪里。
\item 眼睛仍是 $(1.5,1.2,2.5)$，但把粗略上方向改成与 $f$ 平行。\eqref{eq:view} 的哪一次除法会遇到零？这是否表示这样的相机几何上不存在？
\item 同一像素上的颜色若来自三角形三个顶点，屏幕上的重心坐标不能直接拿去平均顶点颜色。原因就在这一讲的那一次除法。试着只对一条线段说明：透视之后，线段中点的像不是两个端点的像的中点。
\end{enumerate}

\appendix
\section{与程序的对应}

程序中的矩阵按列存储。纸上的第 $j$ 列放在数组下标 $4j,4j+1,4j+2,4j+3$。上传时声明不需要转置。

顶点着色器计算的就是 $p_{\mathrm{clip}}=PVM\,p$。片元着色器把顶点颜色原样写出。深度测试和背面剔除是这一讲最后两问的开关，不引入新的公式。

启动时做两道自检。第一道把偏航、俯仰重建的 $V$ 和直接看向原点的 $V$ 逐项比较。第二道是问题「一个角的屏幕坐标」，对八个角各做一遍，姿态取 $R_y(0)$ 和 $R_y(40^\circ)$。任一失败，程序不会进入画图循环。进入循环之后，每帧先用鼠标增量更新 $\psi,\varphi$，再按轴更新 $x$ 和 $z$，然后用重力更新 $y$。任一轴使身体与某个实心方块重叠，就把该轴退回。每个实心方块用平移到 $(i,j,k)$ 的模型矩阵画一次。左键刚按下时，从眼睛沿视线找 \(t\le 5\) 的第一块实心方块，把它改成空气，并立刻重写 \texttt{world.txt}。文件是实心块的完整名单。退出时再写一次。人的位置不在文件里。默认地形由种子 \(20261116\) 生成，不再是手摆的 \(26\) 块。旧版本存档会被这份地形覆盖。

\end{document}
